\section{Individual Component Context}\label{sec:individual-component-context}

The Intent Recognition Engine (IRE) is the first component of the chain
(Component 01). It acts as a domain-specific Natural Language Understanding (NLU)
gateway: it receives natural-language security commands---for example a request
to perform a targeted network discovery scan---and returns a structured,
validated intent contract that downstream components consume without
re-interpreting prose.

The component boundary is deliberately narrow. The IRE is concerned with
understanding and validating \emph{intent}; it does not decide how a task is
scheduled, across which hosts it runs, or how execution failures are recovered.
Those responsibilities belong to the Dynamic Planner (Component 02) and to the
Adaptive Execution and Error Recovery engine (Component 03), which consume the
IRE's output. The handoff is designed to be safe by construction: the output
contract carries an explicit intent type, a resolved target, the recognised
options and any validation findings, and the downstream contract requires that
actions derived for rejected intents must not be executed.

Because the IRE sits at the head of the chain, its reliability bounds the
reliability of the components that follow. A target that is mislabeled, an
intent that is misclassified, or a hallucinated identifier that is accepted would
be inherited by every downstream action. The component therefore couples an
instruction-tuned language model with deterministic validation and guard stages
whose purpose is to contain model error before it can propagate to a planner or
executor. This dissertation evaluates whether that guarded design achieves
sufficiently accurate and error-resilient intent recognition to justify acting on
model-produced security instructions in a controlled, offline setting.